% !TEX program = pdflatex
% !TEX encoding = UTF-8 Unicode

% Plantilla de la clase `scrbook` del paquete KOMA-script para la
% elaboración de un TFG siguiendo las directrices del la comisión del
% Grado en Matemáticas de la Universidad de Granada.

% Francisco Torralbo Torralbo
% miércoles, 29 de abril de 2020

\documentclass{ugrTFG}

% INFORMACIÓN PARA LA VERSIÓN IMPRESA
% Si el documento ha de ser impreso en papel de tamaño a4 pero el tamaño del documento (elegido en \KOMAoptions con la ocpión paper) no es a4 descomentar la línea que carga el paquete `crop` más abajo. El paquete crop se encargará de centrar el documento en un a4 e imprimir unas guías de corte. El procedimiento completo para imprenta sería el siguiente:
% 0. Determinar, según el tipo de encuadernación del documento, el ancho reservado para el proceso de encuadernación (preguntar en la imprenta), es decir, la anchura del área del papel que se pierde durante el proceso de encuadernación. Fijar la varibale BCOR de \KOMAoptions a dicho valor.
% 1. Descomentar la siguiente línea e imprimir una única página con las guías de corte
% 2. Cambiar la opción `cross` por `cam` (o `off`) en el paquete crop y volver a compilar. Imprimir el documento (las guías de corte impresas no inferfieren con el texto).
% 3. Usar la página con las guías impresas en el punto 1 para cortar todas las páginas.

% \usepackage[a4, odd, center, pdflatex, cross]{crop} % Permite imprimir el documento en un a4 (si el tamaño es más pequeño) mostrando unas guías de corte. Útil para imprenta.

% VERSIÓN ELECTRÓNICA PARA TABLETA
% Las opciones siguientes seleccionan un tamaño de impresión similar a una tableta de 9 pulgadas con márgenes estrechos. Útil para producir una versión en pdf para ser leída en una tableta en lugar de impresa.
% Para que la portada quede centrada correctamente hay que editar el
% archivo `portada.tex` y eliminar el entorno `addmargin`

% \KOMAoptions{fontsize=10pt, paper=19.7104cm:14.7828cm, twoside=false, BCOR=0cm, DIV=14}


% Selección del idioma: cargamos por defecto inglés y español (aunque este último es el idioma por defecto para el documento). Cuando queramos cambiar de idioma escribiremos:
% \selectlanguage{english} o \selectlanguage{spanish}




% IMAGEN DE LA PORTADA
% Existen varias opciones para la imagen de fondo de la portada del TFG. Todas ellas tienen en logotipo de la universidad de Granada en la cabecera. Las opciones son las siguientes:
% 1. portada-ugr y portada-ugr-color: diseño con marca de agua basada en el logo de la UGR (en escala de grises y color).
% 2. portada-ugr-sencilla y portada-ugr-sencilla-color: portada únicamente con el logotipo de la UGR en la cabecera.
\usepackage{eso-pic}
\newcommand\BackgroundPic{%
	\put(0,0){%
		\parbox[b][\paperheight]{\paperwidth}{%
			\vfill
			\centering
      % Indicar la imagen de fondo en el siguiente comando
			\includegraphics[width=\paperwidth,height=\paperheight,%
			keepaspectratio]{portada-ugr-sencilla}%
			\vfill
}}}

% \usepackage{listings} % Para la inclusión de trozos de código

% ---------------------------------------------------------------------
% COMANDOS Y ENTORNOS
% ---------------------------------------------------------------------
% Cargamos un archivo externo donde hemos incluido todos los comandos
% propios que vamos a usar en el documento.
% DEFINICIÓN DE COMANDOS Y ENTORNOS

% CONJUNTOS DE NÚMEROS

  \newcommand{\N}{\mathbb{N}}     % Naturales
  \newcommand{\R}{\mathbb{R}}     % Reales
  \newcommand{\Z}{\mathbb{Z}}     % Enteros
  \newcommand{\Q}{\mathbb{Q}}     % Racionales
  \newcommand{\C}{\mathbb{C}}     % Complejos



% --------------------------------------------------------------------
% INFORMACIÓN DEL TFG Y EL AUTOR
% --------------------------------------------------------------------

\newcommand{\miTitulo}{Título del trabajo\xspace}
\newcommand{\miNombre}{Nombre apellidos\xspace}
\newcommand{\miGrado}{Grado en Matemáticas}
\newcommand{\miFacultad}{Facultad de Ciencias}
\newcommand{\miUniversidad}{Universidad de Granada}
% Añadir tantos tutores como sea necesario separando cada uno de ellos
% mediante el comando `\medskip` y una línea en blanco
\newcommand{\miTutor}{
  Nombre del tutor 1 \\ \emph{Departamento del tutor 1} 

  % Añadir tantos tutores como sea necesario. 

  \medskip
  Nombre del tutor 2 \\ \emph{Departamento del tutor 2}
}
\newcommand{\miCurso}{2023-2024\xspace}

% HYPERREFERENCES
% --------------------------------------------------------------------
% Opciones para el paquete hyperref
%----------------------------------

\hypersetup{%
  % hidelinks,            % Enlaces sin color ni borde. El borde no se imprime
  linkbordercolor=0.8 0 0,
  citebordercolor=0 0.8 0,
  citebordercolor=0 0.8 0,
  colorlinks = true,            % Color en texto de los enlaces. Comentar esta línea o cambiar `true` por `false` para imprimir el documento.
  linkcolor = [rgb]{0.5, 0, 0}, % Color de los enlaces internos
  urlcolor = [rgb]{0, 0, 0.5},  % Color de los hipervínculos
  citecolor = [rgb]{0, 0.5, 0}, % Color de las referencias bibliográficas
	pdftitle={\miTitulo},%
	pdfauthor={\textcopyright\ \miNombre, \miFacultad, \miUniversidad},%
  pdfsubject={Trabajo de fin de Grado},%
	pdfkeywords={},%
	pdfcreator={pdfLaTeX},%
}

% Redefinición del estilo para mostrar las referencias cruzadas en la bibliografía.
\renewcommand*{\backref}[1]{}
\renewcommand*{\backrefalt}[4]{{\footnotesize [%
    \ifcase #1 No citado%
    \or Citado en pág.~#2%
    \else Citado en págs. #2%
    \fi%
]}}

% Etiquetas en español para el comando \autoref
\def\chapterautorefname{Capítulo}
\def\appendixautorefname{Apéndice}
\def\sectionautorefname{Sección}
\def\subsectionautorefname{Subsección}
\def\figureautorefname{Fig.}
\def\tableautorefname{Cuadro}

\def\teoremaautorefname{Teorema}
\def\proposicionautorefname{Proposición}
\def\lemaautorefname{Lema}
\def\corolarioautorefname{Corolario}
\def\definicionautorefname{Def.}
\def\observacionautorefname{Observación}
\def\ejemploautorefname{E.j.}

% Pone automáticamente un parántesis para las ecuaciones
\def\equationautorefname~#1\null{Ec.~(#1)\null}


\begin{document}

% --------------------------------------------------------------------
% FRONTMATTER
% -------------------------------------------------------------------
\frontmatter % Desactiva la numeración de capítulos y usa numeración romana para las páginas

% \pagestyle{plain} % No imprime cabeceras

% !TeX root = ../libro.tex
% !TeX encoding = utf8

%*******************************************************
% Titlepage
%*******************************************************
\begin{titlepage}
  \AddToShipoutPicture*{\BackgroundPic}
  \phantomsection 
  \pdfbookmark[1]{Título}{title}

  % Para que el título esté centrado en la página.
  % Los valores numéricos deberán elegirse de acuerdo con el diseño de
  % página (sobre todo si se cambia la opción BCOR o DIV).
  \begin{addmargin}[2.575cm]{0cm}
  \begin{flushleft}
    \Large  
    \hfill\vfil

    \textsf{\miFacultad}
    \vfill\vfill

    {\huge\textsf\miGrado} \vfill


    \textsc{trabajo de fin de grado}

    \begingroup
    \Huge{\miTitulo} \\ \bigskip
    \endgroup

    \vfill\vfill\vfill\vfill

    \textsf{\normalsize{Presentado por:}}\\
    {\Large\textrm{\miNombre}}

    \vfill
    \textsf{Curso académico \miCurso}
  \end{flushleft}  
  \end{addmargin}       

\end{titlepage}   
\cleardoublepage
\endinput
                    
% !TeX root = ../libro.tex
% !TeX encoding = utf8

%*******************************************************
% Little Dirty Titlepage
%*******************************************************

\thispagestyle{empty}

\begin{center}
  \large  

  \vspace*{\stretch{1}}

  \begingroup
  \huge{\miTitulo} \\ \bigskip
  \endgroup

  \textrm{\miNombre}

  \vspace{\stretch{5}}

\end{center}  

\newpage
\thispagestyle{empty}

\hfill

\vfill

\noindent\miNombre \textit{\miTitulo}.

Trabajo de fin de Grado. Curso académico \miCurso.

\begin{minipage}[t]{0.25\textwidth}
  \flushleft
  \textbf{Responsable de tutorización}
\end{minipage}
\begin{minipage}[t]{0.45\textwidth}
  \flushleft
  \miTutor
\end{minipage}
\begin{minipage}[t]{0.30\textwidth}
  \flushright
  \miGrado
  \medskip

  \miFacultad
  \medskip

  \miUniversidad
\end{minipage}
\begin{flushleft}
\end{flushleft}

\endinput
                     
% !TeX root = ../libro.tex
% !TeX encoding = utf8
%
%*******************************************************
% Declaración de originalidad
%*******************************************************

\thispagestyle{empty}

\hfill\vfill

\textsc{Declaración de originalidad}\\\bigskip

D./Dña. \miNombre \\\medskip

Declaro explícitamente que el trabajo presentado como Trabajo de Fin de Grado (TFG), correspondiente al curso académico \miCurso, es original, entendida esta, en el sentido de que no ha utilizado para la elaboración del trabajo fuentes sin citarlas debidamente.
\medskip

En Granada a \today 
\begin{flushleft} 
Fdo: \miNombre 

\end{flushleft}

\vfill

\cleardoublepage
\endinput
   
% !TeX root = ../libro.tex
% !TeX encoding = utf8

%*******************************************************
% Dedication
%*******************************************************
\thispagestyle{empty}
\phantomsection 
\pdfbookmark[1]{Dedicatoria}{Dedicatoria}

\hfill
\vfill

\begin{flushright}
\itshape
Dedicatoria (opcional) \\
Ver archivo \texttt{preliminares/dedicatoria.tex}
\end{flushright}

\vfill

\cleardoublepage
\endinput
                % Opcional
% !TeX root = ../libro.tex
% !TeX encoding = utf8

%*******************************************************
% Table of Contents
%*******************************************************
\phantomsection
\pdfbookmark[0]{\contentsname}{toc}

\setcounter{tocdepth}{2} % <-- 2 includes up to subsections in the ToC
\setcounter{secnumdepth}{3} % <-- 3 numbers up to subsubsections

% \manualmark
% \markboth{\textsc{\contentsname}}{\textsc{\contentsname}}
\tableofcontents 

%*******************************************************
% List of Figures and of the Tables
%*******************************************************

    % *******************************************************
    %  List of Figures
    % *******************************************************    
    \phantomsection 
    \listoffigures

    %*******************************************************
    % List of Tables
    %*******************************************************
    \phantomsection 
    \listoftables
    
    %*******************************************************
    % List of Listings
    % The package \usepackage{listings} is needed
    %*******************************************************      
	  % \phantomsection 
    % \renewcommand{\lstlistlistingname}{Listados de código}
    % \lstlistoflistings 

\cleardoublepage
            
% !TeX root = ../tfg.tex
% !TeX encoding = utf8

%*******************************************************
% Agradecimientos
%*******************************************************

\chapter{Agradecimientos}

Agradecimientos (opcional, ver archivo \texttt{preliminares/agradecimiento.tex}).

\cleardoublepage
\endinput
            % Opcional

% \pagestyle{scrheadings} % A partir de ahora sí imprime cabeceras

% !TeX root = ../libro.tex
% !TeX encoding = utf8
%
%*******************************************************
% Summary
%*******************************************************


\chapter{Summary}

An english summary of the project (around 800 and 1500 words are recommended).

\endinput
                    
% !TeX root = ../libro.tex
% !TeX encoding = utf8
%
%*******************************************************
% Introducción
%*******************************************************

% \manualmark
% \markboth{\textsc{Introducción}}{\textsc{Introducción}} 

\chapter{Introducción}

Introducción del libro

\endinput
               

% --------------------------------------------------------------------
% MAINMATTER
% --------------------------------------------------------------------
\mainmatter % activa la numeración de capítulos, resetea la numeración de las páginas y usa números arábigos

\setpartpreamble[c][0.75\linewidth]{%
	\bigskip % Deja un espacio vertical en la parte superior
  Si el trabajo se divide en diferentes partes es posible incluir al inicio de cada una de ellas un breve resumen que indique el contenido de la misma. Esto es opcional.
}
\part{Primera parte} % Dividir un libro en partes OPCIONAL

% !TeX root = ../libro.tex
% !TeX encoding = utf8

\setchapterpreamble[c][0.75\linewidth]{%
	\sffamily
  Al inicio de cada capítulo puede incluirse un breve resumen. Esto es opcional.
	\par\bigskip
}
\chapter{Primer capítulo}\label{ch:primer-capitulo}

\section{Introducción}
Este documento es una plantilla para la elaboración de un trabajo fin de Grado siguiendo las \href{http://grados.ugr.es/matematicas/pages/infoacademica/tfg/fechaseinstruccionesdefensatfgcurso20172018/!}{directrices} de la comisión de Grado en Matemáticas de la Universidad de Granada que, a fecha de marzo de 2019, son las siguientes:

\begin{itemize}
  \item La  memoria  debe  realizarse  con  un  procesador  de  texto  científico,  preferiblemente (La)TeX.
  \item La portada  debe contener  el  logo  de  la UGR,  incluir  el  título del TFG, el nombre del estudiante y especificar el grado, la facultad y el curso actual.
  \item La contraportada contendrá además el nombre del tutor o tutores.
  \item La memoria debe necesariamente incluir:
    \begin{itemize}
      \item un índice detallado de capítulos y secciones,
      \item un resumen amplio en inglés del trabajo realizado (se recomienda entre 800 y 1500 palabras),
      \item una introducción en la que se describan claramente los objetivos previstos inicialmente en la propuesta de TFG, indicando si han sido o no alcanzados, los antecedentes importantes para el desarrollo, los resultados obtenidos, en su caso y las principales fuentes consultadas,
      \item una bibliografíafinal que incluya todas las referencias utilizadas.
    \end{itemize}
  \item Se recomienda que la extensión de la memoria sea entre 30 y 60 páginas, sin incluir posibles apéndices.
\end{itemize}

Para generar el pdf a partir de la plantilla basta compilar el fichero \texttt{libro.tex}. Es conveniente leer los comentarios contenidos en dicho fichero pues ayudarán a entender mejor como funciona la plantilla. 

La estructura de la plantilla es la siguiente\footnote{Los nombres de las carpetas no se han acentuado para evitar problemas en sistemas con Windows}: 
\begin{itemize}
  \item Carpeta \textbf{preliminares}: contiene los siguientes archivos
    \begin{description}
      \item[\texttt{dedicatoria.tex}] Para la dedicatoria del trabajo (opcional)
      \item[\texttt{agradecimientos.tex}] Para los agradecimientos del trabajo (opcional)
      \item[\texttt{introduccion.tex}] Para la introducción (obligatorio)
      \item[\texttt{summary.tex}] Para el resumen en inglés (obligatorio)
  \end{description}
  El resto de archivos de dicha carpeta no es necesario editarlos pues su contenido se generará automáticamente a partir de los metadatos que agreguemos en \texttt{libro.tex}

  \item Carpeta \textbf{capitulos}: contiene los archivos de los capítulos del TFG. Añadir tantos archivos como sean necesarios. Este capítulo es \texttt{capitulo01.tex}.

  \item Carpeta \textbf{apendices}: Para los apéndices (opcional)
  \item Carpeta \textbf{img}: Para incluir los ficheros de imagen que se usarán en el documento.
  \item Carpeta \textbf{paquetes}: Incluye dos ficheros 
    \begin{description}
      \item[\texttt{hyperref.tex}] para la configuración de hipervínculos al generar el pdf (no es necesario editarlo) 
      \item[\texttt{comandos-entornos.tex}] donde se pueden añadir los comandos y entornos personalizados que precisemos para la elaboración del documento. Contiene algunos ejemplos
    \end{description}
    
  \item Fichero \texttt{library.bib}: Para incluir las referencias bibliográficas en formato \texttt{bibtex}. Son útiles las herramientas \href{https://www.doi2bib.org/}{doi2bib} y \href{https://www.ottobib.com/}{OttoBib} para generar de forma automática el código bibtex de una referencia a partir de su \textsc{doi} o su \textsc{isbn}.

  \item Fichero \texttt{glosario.tex}: Para incluir un glosario en el trabajo (opcional). 

   \item Fichero \texttt{libro.tex}: El documento maestro del TFG que hay que compilar con \LaTeX\ para obtener el pdf. En dicho documento hay que cambiar la \emph{información del título del \textsc{tfg} y el autor así como los tutores}.
\end{itemize}


\section{Elementos del texto}

En esta sección presentaremos diferentes ejemplos de los elementos de texto básico. Conviene consultar el contenido de \texttt{capitulos/capitulo01.tex} para ver cómo se han incluido.

\subsection{Listas}
En \LaTeX\ tenemos disponibles los siguientes tipos de listas:

Listas enumeradas:
\begin{enumerate}
  \item item 1
  \item item 2
  \item item 3
\end{enumerate}

Listas no enumeradas
\begin{itemize}
  \item item 1
  \item item 2
  \item item 3
  \end{itemize}

Listas descriptivas
\begin{description}
  \item[termino1] descripción 1
  \item[termino2] descripción 2
\end{description}
  
\subsection{Tablas y figuras}

En la \autoref{tb:ejemplo-tabla} o la \autoref{fig:logo-ugr} podemos ver\ldots

\begin{table}[htpb]
  \centering
  \begin{tabular}{ccc} \toprule
    \multicolumn{2}{c}{Agrupados} \\ \cmidrule(r){1-2}
    cabecera & cabecera & cabecera          \\ \midrule
    elemento & elemento & elemento          \\ 
    elemento & elemento & elemento          \\ 
    elemento & elemento & elemento          \\ \bottomrule
  \end{tabular}
  \caption{Ejemplo de tabla}
  \label{tb:ejemplo-tabla}
\end{table}

\begin{figure}[htpb]
  \centering
  \includegraphics[width=0.8\textwidth]{logo-ugr}
  \caption{Logotipo de la Universidad de Granada}
  \label{fig:logo-ugr}
\end{figure}

\section{Entornos matemáticos}

\begin{teorema}\label{thm:teorema}
Esto es un ejemplo de teorema.
\end{teorema}

\begin{proposicion}
Ejemplo de proposición
\end{proposicion}

\begin{lema}
Ejemplo de lema
\end{lema}

\begin{corolario}
Ejemplo de corolario
\end{corolario}

\begin{definicion}
Ejemplo de definición
\end{definicion}

\begin{observacion}
Ejemplo de observación
\end{observacion}

Y esto es una referencia al \autoref{thm:teorema}. 


\section{Bibliografía e índice}

Esto es un ejemplo de texto en un capítulo. Incluye varias citas tanto a libros~\cite{Euler1982, Euler1984, Euler1985} como a recursos online~\cite{EulerWiki} (páginas web). Ver el fichero \texttt{library.bib}. 

Además incluye varias entradas al índice alfabético mediante el comando \verb+\index+ \index{Leonard!Euler|textbf}


\endinput

% Añadir tantos capítulos como sea necesario

\cleardoublepage\part{Segunda parte}
% !TeX root = ../libro.tex
% !TeX encoding = utf8

\chapter{Segundo capítulo}

\section{Primera sección}

\endinput
%--------------------------------------------------------------------
% FIN DEL CAPÍTULO. 
%--------------------------------------------------------------------


% --------------------------------------------------------------------
% APPENDIX: Opcional
% --------------------------------------------------------------------

\appendix % Reinicia la numeración de los capítulos y usa letras para numerarlos
\pdfbookmark[-1]{Apéndices}{appendix} % Alternativamente podemos agrupar los apéndices con un nuevo \part{Apéndices}

% !TeX root = ../libro.tex
% !TeX encoding = utf8

\chapter{Primer apéndice}\label{ap:apendice1}

Los apéndices son opcionales.

Archivo: \texttt{apendices/apendice01.tex}

\endinput
%------------------------------------------------------------------------------------
% FIN DEL APÉNDICE. 
%------------------------------------------------------------------------------------

% Añadir tantos apéndices como sea necesario 

% --------------------------------------------------------------------
% GLOSARIO: Opcional
% --------------------------------------------------------------------

% !TeX root = ../libro.tex
% !TeX encoding = utf8

\chapter*{Glosario}
\addcontentsline{toc}{chapter}{Glosario} % Añade el glosario a la tabla de contenidos

La inclusión de un glosario es opcional.

Archivo: \texttt{glosario.tex}

\begin{description} 
  \item[$\mathbb{R}$] Conjunto de números reales.

  \item[$\mathbb{C}$] Conjunto de números complejos.

  \item[$\mathbb{Z}$] Conjunto de números enteros.
\end{description}
\endinput
 

% -------------------------------------------------------------------
% BACKMATTER
% -------------------------------------------------------------------

\backmatter % Desactiva la numeración de los capítulos
\pdfbookmark[-1]{Referencias}{BM-Referencias}

% BIBLIOGRAFÍA
%-------------------------------------------------------------------

\setbibpreamble{Las referencias se listan por orden alfabético. Aquellas referencias con más de un autor están ordenadas de acuerdo con el primer autor.\par\bigskip}
\bibliographystyle{alpha} 
\begin{small} % Normalmente la bibliografía se imprime en un tamaño de letra más pequeño.
\bibliography{library.bib}
\end{small}


\end{document}
